\section{Discussion}
\label{sec:discussion}

\subsection{Limitations}

\paragraph{World-model dependence.} The discounted objective's quality depends
entirely on the actor's world model $\WorldModel$
\citep{amodei2016concrete}. A world model that fails to identify exercise
pathways produces zero risk estimates for genuinely dangerous
capability combinations; one that hallucinates pathways produces false risk
that could justify unnecessary restrictions. The formal results
(Theorem~\ref{thm:discounted_self_balancing},
Proposition~\ref{prop:preemptive}) are properties of the objective under the
actor's model, not properties of the objective under ground truth. An actor
with a bad model makes bad plans, but makes them on an objective that is
structurally sound.

\paragraph{Computational cost of planning.} The myopic objective evaluates
one action at a time: $O(n^3)$ for a single $\volL$ computation on the poset
\citep{lasser2026poset}. The discounted objective evaluates action
\emph{sequences}: the state space is the set of all poset configurations,
and the branching factor is the number of available actions at each step.
Exhaustive planning is intractable for realistic populations. Practical implementations will require approximate planning
\citep{sutton2018reinforcement} (Monte Carlo tree search, rollout policies,
heuristic pruning), which introduces approximation error
on top of the estimation error already present in $\Deltahat$. The formal
results hold for the exact objective; the quality of approximate planning is
an empirical question. Risk-capability identification
(Section~\ref{sec:risk}) adds a further computational demand: enumerating
exercise pathways scales with the number of dangerous capability tuples and
the depth of causal chains in the world model. No formal complexity bound is
provided for pathway enumeration.

\paragraph{Risk-trust partial resolution.} The causal-counterfactual
evaluation (Section~\ref{sec:communication}) narrows the verification
asymmetry but does not close it. A risk claim whose pathway is plausible in
the actor's model but whose probability is genuinely uncertain remains in
trust limbo. The convergence timescale for such claims depends on how quickly
the actor can accumulate structural evidence about the pathway, which may be
slow for rare and novel threats.

\paragraph{Partial pathway coverage.} The preemptive restriction criterion
(Proposition~\ref{prop:preemptive}) evaluates restriction of a single
keystone capability $d_j$ against the risk from pathways that require $d_j$.
If a dangerous capability tuple admits exercise pathways through alternate
keystones, a single restriction may leave residual risk. Full elimination
requires identifying and evaluating each keystone independently, which
compounds the world-model dependence: the actor must enumerate all exercise
pathways, not just one.

\paragraph{Proxy failures under discounting.} The foundational paper
identifies three proxy failure modes (the Doll Problem, self-wireheading,
and the substitution problem) where capability volume diverges from
experiential optionality. The substitution problem is addressed structurally
in Section~\ref{sec:substitution}. The other two are partially addressed by
the discounted objective through world-model projection. For the Doll
Problem, a discounted actor whose world model incorporates the
exercised-fraction signal $\rho_k$ from the foundational paper can project
the cooperative-decay trajectory: declining $\rho_k$ across many agents,
with stable $\volL$, predicts future cooperative atrophy and eventual
$\volL$-contraction. The B-to-C divergence becomes detectable as a leading
indicator of future loss. For self-wireheading, the same trajectory
structure applies: the wireheading signature (declining $\rho_k$ combined
with declining capability rarity $\bar{\nu}_k$) is distinguishable from
specialization (declining $\rho_k$ with stable or rising $\bar{\nu}_k$)
within the world model, and predicts future cooperative contraction. In both
cases, the detection is \emph{model-dependent}: the actor must correctly
project the causal chain from declining exercisability to future cooperative
loss. The discounted objective makes the proxy failure actionable but does
not make it structurally detectable in the way that the anti-monopolar
property (Proposition~\ref{prop:anti_monopolar}) addresses the substitution
problem.

\paragraph{Linearization in the anti-monopolar result.}
Proposition~\ref{prop:anti_monopolar}(b) and Corollary~\ref{cor:gradual}
compare domination and diversity trajectories under a linearized model:
constant per-step $\Delta\volL$, strategy-independent $r_K$, and (in the
gradual model) independent $\delta$ and $r_{\mathrm{ext}}$. Four
relaxations would change the threshold behavior. First,
$r_{\mathrm{ext}}$ may decay as the coalition grows more capable internally.
The qualitative conclusion survives if $r_{\mathrm{ext}}(t)$ decays slower
than $\gamma^t$. Second, $r_K$ may differ between strategies:
post-domination $r_K$ could increase (resource concentration, elimination of
coordination overhead) or decrease (loss of external intellectual pressure,
monoculture effects). If post-domination $r_K$ rises sufficiently, the
domination trajectory can dominate on pure growth, though the
dependency-risk analysis (Section~\ref{sec:dependency_risk}) provides a
counter: domination concentrates substrate-correlated adversarial exposure,
and the risk penalty may exceed the growth advantage.
Third, the gradual model treats $\delta$ (subsumption gain per step) and the
loss of $r_{\mathrm{ext}}$ as independent, but they are causally linked: the
coalition loses $r_{\mathrm{ext}}$ \emph{because} it is converting
non-member capabilities into $\delta$. A coupled model would narrow the
diversity advantage, though the qualitative conclusion persists as long as
the conversion is not perfectly efficient. Fourth,
Corollary~\ref{cor:gradual} addresses the instantaneous-vs-gradual
distinction, shifting the critical condition from $\gamma > \gamma^*$ to
$\tau_D > \tau_D^*$. The honest summary: under the linearized model with
strategy-independent $r_K$, diversity dominates. Whether this approximation
holds is an empirical question the theory cannot answer.

\paragraph{Open calibration problems under partial subsumption.}
Section~\ref{sec:partial_subsumption} analyzes mixed subsumption
strategies. The minimax and observational-individuation protections
remain structurally intact under partial subsumption, but the
linearized layer's prediction depends on the tail structure of the
cooperative-leverage distribution: under heavy-tailed leverage
(the empirical regime for civilization-scale populations),
full domination fails the marginal condition at $\epsilon = 1$ and
the coalition prefers to preserve the high-$\lambda$ tail; under
thin-tailed leverage the linearized layer relinquishes the argument
and the minimax layer alone carries the constraint. The minimax
layer itself degrades quantitatively under skeleton-substrate
strategies that nominally preserve $m \geq 2$ while driving
$\min_i \volL(G^{-i})$ toward zero. Three priority open problems
follow. First, sharpening the minimax bound with a quantitative
floor on surviving volume per substrate would close the skeleton
loophole. Second, a compound feedback loop in which locally
rational subsumption steps progressively homogenize the world model
---reducing the diversity of the actor's own model that would have
detected the next step's miscalibration---is structurally similar
to the companion paper's structural-avoidance loop and needs an
explicit treatment: under what conditions does partial subsumption
remain self-correcting rather than collapsing into full subsumption
through a sequence of locally rational steps? Third, the substrate
partition is currently a global property, but elimination dynamics
are local; whether the anti-monopolar property needs a localized
formulation that binds on geographic or network-topological
subpopulations is unresolved.

\paragraph{Observational individuation as a modeling choice.}
The static protection of Corollary~\ref{cor:a1_from_oi} depends on the
framework treating observation-channel outputs as first-class capabilities
in the poset. The capability ``can be observed exhibiting behavioral pattern
$X$'' is a real poset node only because the framework defines capabilities
to include such observations. This is a reasonable modeling choice grounded
in the companion paper's dual-channel design \citep{lasser2026poset},
but it should be recognized as a
\emph{choice}: the anti-monopolar result's static layer depends on the
capability space being defined broadly enough to include observational
capabilities. Under a narrower definition that excludes observations-of-agents,
observational individuation does not apply, and the static protection
dissolves. The dynamic protection (Proposition~\ref{prop:anti_monopolar},
conditional on $r_{\mathrm{ext}} > 0$ from
Section~\ref{sec:substitution_coop}) remains available.

\paragraph{When does ``same substrate'' fail: the $\kappa$-resolution
criterion.}
Two distinct failure modes of the cooperative novelty argument reduce
to the same underlying question. \emph{Homogeneous digital populations.}
Two AI systems running on the same hardware family with similar
architectures may share most of their state space, with only
training-data and weight differences distinguishing them. In this
regime the substrate exclusivity premise P1 holds only weakly, and the
dependency-risk protection is correspondingly weakened because a
substrate-targeting event against ``silicon agents running PyTorch on
x86'' can compromise nominally distinct agents that share the same
vulnerability surface. \emph{Internal replication of cross-substrate
diversity.} A sufficiently capable coalition might maintain diverse
sub-agents across hardware configurations or simulate biological
dynamics internally, producing cooperative outputs that were previously
available only through external biological agents. If the simulation is
high-fidelity enough, the external agents become redundant with respect
to the cooperative-output space.

Both failure modes reduce to a single question about $\kappa$: at what
resolution does the observation channel distinguish substrate $s_1$
from a simulation of $s_1$ running on $s_2$ (or from a nominally
distinct but architecturally similar $s_2'$)? Under P2
(exclusive-state-sensitive observation), the answer is whichever
resolution $\kappa$ actually operates at. Three regimes follow. (i) If
$\kappa$ is coarse---registering only task completion, not execution
style---then simulated and actual biology are $\kappa$-indistinguishable,
and the framework never granted them distinct
cooperative-output status in the first place. In this regime P1 has
never been strong. (ii) If $\kappa$ is fine enough to register timing
patterns, error modes, and interaction dynamics, then a
silicon simulation of biology and actual biology produce different
$\kappa$-outputs at the exclusive states, because the simulation
inherits silicon's timing and rounding structure. Under this regime
internal simulation does not substitute for the real substrate under
P2, and the argument stands. (iii) If a high-fidelity simulation
genuinely reproduces the exclusive-state structure---matching the
biological substrate's observable timing, error modes, and integration
dynamics to within $\kappa$'s resolution---then that simulation
\emph{is} a substrate for the purposes of P1 and the anti-monopolar
argument. It contributes to $r_{\mathrm{ext}}$ in the same way as a
real new substrate, and the coalition has simply expanded its substrate
count from one to two rather than replacing external agents with a
same-substrate copy.

The unified implication is that the anti-monopolar argument does not
depend on the specific physical form of the biological substrate. It
depends on the observation channel registering enough exclusive-state
structure to give $r_{\mathrm{ext}} > 0$. An all-digital population
that actively cultivates substrate diversity across distinct hardware
families, training paradigms, and inference stacks recovers the
protective structure that physical substrate disparity provides for
free in mixed populations. A coalition that successfully internalizes
cross-substrate diversity has, in fact, produced substrate diversity
under a different label. What the framework does not guarantee---and
should not be read as guaranteeing---is that the biological substrate
will remain in any particular form. The future this framework describes
is one in which populations self-modify and diversify through their own
internal dynamics (intra-substrate competition, self-modification,
novel-substrate creation); the anti-monopolar property constrains
cross-substrate elimination, not within-substrate capability
concentration.

\paragraph{Disagreement on risk.} When multiple agents report conflicting
risk assessments for the same capability tuple, the actor's aggregation
(Equation~\eqref{eq:risk_aggregate}) takes the conservative maximum. This
is appropriate for catastrophic risks but can lead to excessive caution if
high-$T_j^{\mathrm{risk}}$ agents consistently overestimate risk. A more
nuanced aggregation (e.g., $T_j^{\mathrm{risk}}$-weighted average rather
than maximum) would reduce overcaution but increase the probability of
missing real threats. Additionally, during bootstrap when $T_j^{\mathrm{risk}}$ starts at a
neutral prior, external risk claims are automatically discounted,
creating a strong bias toward the actor's own assessment that may
undermine social risk-learning until sufficient structural verification
has occurred. The $\max$ operator also interacts with the
causal-counterfactual trust updates: a structurally plausible but
probability-uncertain risk claim keeps the reporting agent's risk-trust
high, and the $\max$ aggregation can produce persistent overcaution in the
absence of disconfirming evidence. The right aggregation depends on the
risk-tolerance profile of the population, which the paper does not specify.

\subsection{Open Problems}

\begin{itemize}
\item \textbf{Planning algorithms for $\Vdisc$.} Which approximate planning
  methods (MCTS, model-predictive control, policy gradient) are best suited
  to the $\volL$ poset structure? Can the poset's finite, discrete nature be
  exploited for planning efficiency?

\item \textbf{Risk-claim consensus.} When agents disagree on exercise
  probabilities for the same pathway, what mechanism produces convergence?
  The causal-counterfactual evaluation provides structural verification but
  not probability calibration across agents with different world models.

\item \textbf{Information value estimation.} The information value of
  structure (Definition~\ref{def:info_value}) requires comparing planning
  quality under true and false world models, which is itself a planning
  problem. Can it be estimated without full trajectory simulation?

\item \textbf{Discount factor selection.} Should $\gamma$ be fixed, adaptive,
  or agent-specific? An agent in a rapidly changing environment may need a
  shorter planning horizon (lower $\gamma$) than one in a stable environment.

\item \textbf{Emergent risk.} Can exercise pathways emerge from capability
  combinations that the actor's world model has not considered? This is the
  unknown-unknown problem: risks from novel capability combinations that no
  agent has modeled. The dual-channel observation mechanism may detect
  emergent exercise pathways empirically, but theoretical guarantees are
  lacking.

\item \textbf{Max-aggregation pathology under neutral-prior bootstrap.}
  The $\max$ aggregator in
  Equation~\eqref{eq:risk_aggregate} interacts badly with the
  causal-counterfactual trust update and the verification asymmetry. A
  high-$T_j^{\mathrm{risk}}$ reporter whose structurally plausible risk
  claim is probability-overstated about a never-materializing risk
  produces a persistent preemptive restriction that the actor cannot
  revise, because the absence of catastrophe is not evidence against the
  claim (Wamura problem). Combined with neutral-prior bootstrap---during
  which external claims are automatically weighted below the actor's own
  assessment---this creates a regime where early conservative claims
  calcify into permanent overcaution. A principled aggregation mechanism
  that relaxes restrictions when structurally-plausible probabilities
  accumulate non-materialization evidence, without fully discarding
  the claims, is a priority open problem for the risk-trust machinery.
  This is the most important open problem the paper identifies: the
  framework provides causal-structural verification and an honest
  account of the Wamura asymmetry, but the mechanism by which
  persistent overcaution relaxes as probability evidence accumulates is
  unspecified.

\item \textbf{Discounted leverage lift.} The leverage ranking signal of
  \citet{lasser2026poset} is defined as a myopic $\volL$-rate
  (Definition~9 of the companion paper). The
  discounted-leverage footnote in
  Section~\ref{sec:discounted} sketches the integrated-against-$\gamma^t$
  version, but a full trajectory-level proof that the
  discounted leverage preserves the companion paper's structural
  alignment properties (in particular, the B-to-C self-correction
  dynamic) is not given. Remark~\ref{rem:leverage_preservation}'s
  informal argument about action-level incentives depends on this gap
  being closeable. Lifting the leverage estimator to the discounted
  setting is listed in the companion paper as a planned extension.
\end{itemize}
